\documentclass[10pt]{amsart}
\usepackage[T1]{fontenc}
\usepackage{amsmath,amssymb,amsthm}
\usepackage[margin=1in]{geometry}
\usepackage{hyperref}
\hypersetup{pdftitle={Explicit upper and lower bounds for distinct consecutive sums in permutations},pdfauthor={Haoyu Chen},colorlinks=true,linkcolor=blue,citecolor=blue,urlcolor=blue}
\allowdisplaybreaks
\newtheorem{theorem}{Theorem}
\newtheorem{lemma}{Lemma}[section]
\theoremstyle{remark}
\newtheorem{remark}[lemma]{Remark}
\newcommand{\Fp}{\mathbb F_p}
\newcommand{\E}{\mathbb E}
\newcommand{\ind}{\mathbf 1}
\newcommand{\normA}[1]{\lVert #1\rVert_A}
\newcommand{\pos}[1]{(#1)_+}
\newcommand{\Sset}{\mathcal S}
\newcommand{\Uset}{\mathcal U}
\newcommand{\Cset}{\mathcal C}
\newcommand{\Bset}{\mathcal B}
\newcommand{\Aset}{\mathcal A}
\title[Distinct consecutive sums in permutations]{Explicit upper and lower bounds for distinct consecutive sums in permutations}
\author{Haoyu Chen}
\address{Independent researcher}
\date{12 October 2026}
\subjclass[2020]{11B75, 05A05, 11T23}
\keywords{Consecutive sums, permutations, explicit bounds, exponential sums}
\begin{document}
\begin{abstract}
Let $f(n)$ be the largest number of distinct nonempty consecutive sums in
a permutation of $\{1,\ldots,n\}$. We prove
$f(n)\le (4435834/10^7)n^2+1100n$ for every $n\ge1$, and
$f(n)\ge0.309799n^2$ for every $n\ge2^{4096}$.
The upper bound uses coupled rearrangement inequalities and an exact
integer certificate. The lower bound uses constant-sum triples ordered
by affine images of primitive-root powers, with the zero field label
retained, and Fourier completion with restricted translations.
At length $3p$, for every prime $p\ge2^{16}$, the construction gives
$(803/288)p^2-50000p^{7/4}(\log p+3)^8$ distinct sums; this expression
becomes positive only at about $p=2^{304.22}$, and gives the coefficient
$0.309799$ for $p\ge2^{416}$. Dusart's explicit prime-interval theorem
supplies the extension to all lengths. Both limiting coefficients
improve those of Konieczny, the strongest coefficients we located in a
bounded literature search completed on 12 October 2026. The search was
not exhaustive and does not establish priority. Neither the limiting
value of $f(n)/n^2$ nor its existence is determined.
\end{abstract}
\maketitle

\section{Introduction and statements}\label{sec:intro}
For a permutation $a=(a_1,\ldots,a_n)$ of $\{1,\ldots,n\}$, put
\[
 \Sset(a)=\{a_{i+1}+\cdots+a_j:0\le i<j\le n\},
 \qquad f(n)=\max_a |\Sset(a)|.
\]
Intervals are linear and nonempty, and sums are counted by distinct
values, without multiplicity. This extremal quantity occurs in
Erd\H{o}s Problem~34 \cite{Bloom}. All logarithms below are natural.
Every error in our bounds is explicit.

\begin{theorem}[Upper bound]\label{thm:upper}
For every integer $n\ge1$,
\[
 f(n)\le\frac{4435834}{10^7}n^2+1100n.
\]
Consequently $f(n)\le0.4436n^2$ for every integer $n\ge66\,265\,061$.
\end{theorem}

\begin{theorem}[Prime lengths]\label{thm:prime}
Let $p\ge2^{16}$ be prime, $k=(p-1)/2$, $m=\lfloor p^{3/4}\rfloor$,
and $I_p=\{k-m,\ldots,k+m\}$. The explicit family $a(p,c,t)$ of
Section~\ref{sec:construction}, with $c\in\Fp^*$, $t\in I_p$, and
the label $x_{p-1}=0$ retained, satisfies
\begin{equation}\label{eq:mainlower}
 \max_{c\in\Fp^*,\ t\in I_p}|\Sset(a(p,c,t))|
 \ge \frac{803}{288}p^2-50000p^{7/4}(\log p+3)^8.
\end{equation}
In fact the average over this parameter set has the same lower bound.
For every prime $p\ge2^{416}$, a member of this family has at least
$0.309799(3p)^2$ distinct sums.
\end{theorem}

The threshold $2^{16}$ in Theorem~\ref{thm:prime} is a technical side
condition, not an onset of numerical content. For an exact description,
let $u_*>29$ be the unique solution of
\begin{equation}\label{eq:positivity}
 50000e^{-u_*/4}(u_*+3)^8=803/288.
\end{equation}
Within the theorem's range the right side of \eqref{eq:mainlower} is
positive exactly when $p>e^{u_*}$. Numerically
$u_*/\log2\approx304.220$; a directed integer-series check encloses
it strictly between $304.219$ and $304.221$. These numbers describe
the real threshold, not a particular prime. Section~\ref{sec:onsets} proves the exact
criterion and the stated effective onsets. At all experimentally
accessible lengths the error-bearing bound is vacuous.

\begin{theorem}[All lengths]\label{thm:all}
For every integer $n\ge2^{4096}$,
\[
 f(n)\ge0.309799n^2.
\]
Consequently
\[
 0.309799\le\liminf_{n\to\infty}\frac{f(n)}{n^2}
 \le\limsup_{n\to\infty}\frac{f(n)}{n^2}\le0.4435834.
\]
\end{theorem}

The proof of Theorem~\ref{thm:all} uses precisely the following
published input, with the indicated endpoints:
\begin{equation}\tag{P}\label{eq:P}
 x\ge3275\quad\Longrightarrow\quad
 \text{there is a prime }p\text{ with }
 x<p\le x\left(1+\frac1{2(\log x)^2}\right).
\end{equation}
This is Dusart's 1998 thesis, Th\'eor\`eme~1.9, printed p.~35
\cite{Dusart}. The prime-length result does not use this input.
The maximizing parameters in Theorem~\ref{thm:prime} can be selected
by finite search with an exact tie rule; no efficient algorithm at
these large onsets is claimed.

\section{Related work}\label{sec:related}
Konieczny \cite[Theorem~1.2, p.~415]{Konieczny} proved
\[
 \frac32-\frac2{\sqrt e}\le\liminf_{n\to\infty}\frac{f(n)}{n^2}
 \le\limsup_{n\to\infty}\frac{f(n)}{n^2}
 \le\frac14+\frac\pi{16}.
\]
The endpoints are $0.286938680574\ldots$ and $0.446349540849\ldots$.
The preprint first appeared in 2015; the journal version appeared in
2021. Our bounds improve the strongest coefficients we located in a
bounded literature search completed on 12 October 2026, namely these
coefficients of Konieczny. The search was not exhaustive and does not
establish priority. It covered the problem's public records, relevant
title and abstract searches, forward-citation records and selected
public mathematical repositories. Such a search can miss unindexed
work or work using different terminology.

Hegyv\'ari's earlier quadratic-residue construction
\cite[Theorem~1, pp.~193--194]{Hegyv} gives coefficient $1/18$ after
extension to a full permutation, as explained in
\cite[Section~1.5, p.~417]{Konieczny}. Its short sequence has the form
$a_{i+1}=2p+r_{i+1}-r_i$, where $r_i=[i^2]_p$, and prefix sums
$2pi+r_i$. Algebraic labels and distinct-difference methods are thus
historical ingredients. The fixed-lag injectivity of primitive-root
powers is also classical Welch--Costas machinery; see
\cite[Theorem~7.29]{Drakakis} and the construction attributed there to
Golomb \cite{Golomb}. We claim no originality for that ingredient.

Konieczny's lower construction randomizes complementary pairs
\cite[Section~4.2, Proposition~4.2]{Konieczny}. The construction here
uses complementary triples, a restricted translation average and
explicit restoration of overlaps already covered by a base set.
The upper reduction in \cite[Section~5.1, (97)]{Konieczny} counts low
values and high-sum intervals. We retain this reduction and add
coupled information about two high intervals and lower-threshold
intervals. Sharpness of the auxiliary high-interval count is different
from sharpness for distinct sums, as already emphasized in
\cite[Remark~5.2]{Konieczny}. The original assertion $f(n)=o(n^2)$ was
already disproved; the present paper sharpens coefficients and does
not claim a solution of the remaining extremal questions.

\section{The upper bound}\label{sec:upper}
\subsection{Rearrangement and the baseline deficit}
Fix $a$, use half-open position intervals, and let $s(X)$ denote the
sum on a set of positions $X$. Set
\[
 T_n=n(n+1)/2,\quad U_n(r)=r(2n-r+1)/2,\quad V(r)=r(r+1)/2.
\]
Every $r$ positions have sum in $[V(r),U_n(r)]$. Weighted
rearrangement follows by swapping an inverted pair: if weights
$c_i>c_j$ receive entries $a_i<a_j$, the swap gains
$(c_i-c_j)(a_j-a_i)>0$. This applies also to negative weights.
Two length-$k$ intervals at integer start separation $\delta\le k$
therefore satisfy
\begin{equation}\label{eq:pair}
 s(\Aset)+s(\Bset)\le U_n(k-\delta)+U_n(k+\delta)
 =2nk-k^2-\delta^2+k.
\end{equation}

Let $H_k$ count length-$k$ intervals with sum $>T_n/2$, and put
\[
 x_*=1-1/\sqrt2,\qquad b(x)=\sqrt{1/2-(1-x)^2},
 \qquad
 g_0(x)=\begin{cases}
 0,&x<x_*,\\ b(x),&x_*\le x\le1/2,\\1-x,&1/2\le x\le1.
 \end{cases}
\]
Two such intervals cannot be disjoint. For $k\le n/2$, their extreme
starts have separation $\delta$ satisfying
\[
 \delta^2<n^2 b(k/n)^2+k-n/2\le n^2b(k/n)^2.
\]
Applying the same calculation to a single interval twice excludes
$k/n<x_*$. A set of starts of span $\delta$ has at most $\delta+1$
members, so
\[
 e_k=g_0(k/n)-(H_k-1)_+/n\ge0.
\]
The function $g_0$ has variation $1$ and integral $\pi/16$.
Indeed, after putting $r^2=1/2$ and $u=1-x$, the triangle and circular
segment have combined area $\pi r^2/8$; integration of the circular
part uses $u=r\sin\theta$. For variation $V_g$, the normalized
right-endpoint sum differs from the integral by at most $V_g/n$:
on each cell bound the error by its width times its variation, and add.
Consequently
\[
 \sum_k H_k\le(\pi/16)n^2+2n-n\sum_ke_k.
\]
There are at most $T_n/2$ integer values at or below $T_n/2$; count
higher values by representing intervals. With $B_0=1/4+\pi/16$,
\begin{equation}\label{eq:baseline}
 |\Sset(a)|\le B_0n^2+\frac94n-n\sum_ke_k.
\end{equation}

\subsection{Anchors and general complement inequalities}
Suppose a bin $[\alpha,\beta)\subset(x_*,1/2)$ contains a length $k/n$
with $e_k\le v<b(\alpha)$. Then $H_k\ge2$. Choose its extreme high
intervals $\Aset=[s,s+k)$ and $\Bset=[t,t+k)$, and set
\[
 x=k/n,\quad d=(t-s)/n,\quad I=\Aset\cap\Bset,\quad
 \Uset=\Aset\cup\Bset,\quad c=1-x-d.
\]
Here $|I|/n=x-d$, $|\Uset^c|/n=c$, and
\[
 b(x)-v\le d<b(x),\qquad d\ge b(\alpha)-v>0.
\]
The positive rank deficit is
\begin{equation}\label{eq:rankdef}
 E_u=U_n(|I|)+U_n(|\Uset|)-T_n
 =n^2\{b(x)^2-d^2\}+k-n/2
 \le n^2\{b(x)^2-d^2\}.
\end{equation}
Also $s(I)>s(\Uset^c)\ge V(nc)$, since the anchor sums total $>T_n$.

Choose rational $\varepsilon$ satisfying
\begin{equation}\label{eq:epscond}
 0<\varepsilon<1/16,\qquad
 4\varepsilon\le(1-\beta-b(\beta))^2,\qquad
 \varepsilon\le2(1/2-\beta)^2.
\end{equation}
Here $1-\beta-b(\beta)>0$: for $\beta<1/2$, $(1-\beta)^2>1/4$
and $b(\beta)^2=1/2-(1-\beta)^2<(1-\beta)^2$.
Put $Q=(1/4-\varepsilon)n^2+n/2$. Suppose two disjoint length-$j$
intervals $L,R$ have sums $>Q$. Let $G$ be their intervening gap,
$\Cset=(L\cup R)^c$, and
\[
 \ell=j/n=1/2-y,\qquad h=2y,\qquad z=|G|/n.
\]
Since $|\Cset|=nh$,
\begin{equation}\label{eq:complementmass}
 s(\Cset)<T_n-2Q=2\varepsilon n^2-n/2,
 \qquad 0\le y<\sqrt\varepsilon.
\end{equation}
The latter follows from $s(\Cset)\ge n^2h^2/2$.
The function $1-x-b(x)$ decreases, so \eqref{eq:epscond} gives $h\le c$ and
\begin{equation}\label{eq:coremass}
 s(I)>V(nc)\ge n^2c^2/2\ge2\varepsilon n^2>s(\Cset).
\end{equation}
Thus $G$ cannot contain $I$.

Use the integer counts
$Z=|\Uset^c|$, $m_C=|\Cset|$, $J=|\Cset\cap\Uset|$,
$K_C=|\Cset\cap I|$, and $W_a=2nd$.
Rearrangement with weights $\ind_{\Aset}+\ind_{\Bset}-\ind_{\Cset}$ gives
\begin{align}
 s(\Cset)&>H_1-E_u,\label{eq:H1}\\
 H_1&=V(m_C-J)+Z(J-K_C)+V(J-K_C)
 +(Z+W_a)K_C+V(K_C).\notag
\end{align}
Indeed its weighted maximum is
$U_n(|I|-K_C)+U_n(|\Uset|-(J-K_C))-V(m_C-J)$.
Subtracting from $U_n(|I|)+U_n(|\Uset|)$ proves the formula.
Weights $\ind_{\Aset}+\ind_{\Bset}-2\ind_{\Cset}$ similarly give
\begin{align}
 2s(\Cset)&>H_2-E_u,\label{eq:H2}\\
 H_2&=V(m_C-J)+V(m_C-K_C)+V(Z+J)+V(Z+W_a+K_C)
       -V(Z)-V(Z+W_a).\notag
\end{align}
The weighted maximum here is
$U_n(|I|-K_C)+U_n(|\Uset|-J)-V(m_C-K_C)-V(m_C-J)$.
Listing the weights $2,1,0,-1,-2$ proves both identities, including
the finite linear terms, without any assumption $j\ge k$.

Put $u=J/n$, $w=K_C/n$. Discarding nonnegative linear terms in
\eqref{eq:H1} and using \eqref{eq:rankdef} yields
\begin{equation}\label{eq:quad1}
 \frac{s(\Cset)}{n^2}>
 \frac{(h-u)^2}{2}+c(u-w)+\frac{(u-w)^2}{2}
 +(c+2d)w+\frac{w^2}{2}-b(x)^2+d^2.
\end{equation}
For completeness, the corresponding expression from \eqref{eq:H2} is
\begin{equation}\label{eq:quad2}
 \frac{2s(\Cset)}{n^2}>
 \frac{(h-u)^2}{2}+\frac{(h-w)^2}{2}+cu+\frac{u^2}{2}
 +(c+2d)w+\frac{w^2}{2}-b(x)^2+d^2.
\end{equation}
This uses $V(Z+J)-V(Z)=ZJ+V(J)$ and its analogue with $K_C$.
In \eqref{eq:quad1}, the derivatives before the constant term are
$c-h+2u-w\ge0$ and $2d-u+2w\ge0$ on the feasible domain
$w\le u$, $w\ge\max(0,u-2d)$, since $h\le c$.
For the second derivative, when $u\le2d$ use $w\ge0$; when $u>2d$
use $w\ge u-2d$. Lower bounds $u\ge z$, $w\ge(z-\eta)_+$,
$0\le\eta\le d$, can therefore be substituted together: minimize
first over $w\ge\max((z-\eta)_+,u-2d,0)$ and then over $u\ge z$.
Along the moving edge $w=u-2d$ the derivative is $c-h+2u\ge0$;
at $u=z$ the maximum is $(z-\eta)_+$ because $\eta\le d\le2d$.
The derivatives of \eqref{eq:quad2} are $c-h+2u$ and
$c+2d-h+2w$, also nonnegative.
If $G\subset I$, substitute $u=w=z$ and use
\eqref{eq:complementmass} to obtain
\begin{align}
 z^2+(1-x+d-h)z+h^2/2-b(x)^2+d^2&<2\varepsilon,\label{eq:full1}\\
 z^2+(1-x-h)z+h^2/2-\{b(x)^2-d^2\}/2&<2\varepsilon.\label{eq:full2}
\end{align}

\subsection{Three bounds on a protruding wing}
If $G$ extends to the left of $I$, then $L$ is disjoint from $\Bset$;
the right case is its reflection. Define
\[
 D_a=t-s,\quad R_0=n-k-j,\quad
 V_*= (L\cup\Bset)^c,\quad P=|V_*\cap(\Aset\setminus\Bset)|.
\]
Thus $0\le P\le D_a$, $P\le R_0$, and
\begin{equation}\label{eq:smallV}
 s(V_*)<T_n/2-Q=\varepsilon n^2-n/4.
\end{equation}
The weights $\ind_{\Aset}+\ind_{\Bset}-\ind_{V_*}$ give
\begin{equation}\label{eq:wingseed}
 s(V_*)>P^2+(|\Uset^c|-R_0)P+V(R_0)-E_u.
\end{equation}
The expression before subtracting $E_u$ is exactly
$V(R_0-P)+|\Uset^c|P+V(P)$.
Two further rearrangements give
\begin{align}
 s(\Aset)+s(\Bset)+s(L)&\le U_n(k-P)+U_n(k+j+P),\label{eq:triple1}\\
 s(\Aset)+s(\Bset)+2s(L)&\le
 U_n(D_a-P)+U_n(k-D_a+j)+U_n(k+j+P).\label{eq:triple2}
\end{align}
The displayed arguments are the superlevel sizes of the weights.
Their largest size is at most $n$ because $P\le R_0$.
Normalize $\eta=P/n$, $r=1-x-\ell$, and $U_0(v)=v-v^2/2$.
The strict thresholds and \eqref{eq:rankdef} imply
\begin{align}
 \eta^2+(\ell-d)\eta+r^2/2-b(x)^2+d^2&<\varepsilon,\label{eq:seednorm}\\
 \eta^2+\ell\eta&<U_0(x)+U_0(x+\ell)-3/4+\varepsilon,\label{eq:triple1norm}\\
 \eta^2+(x+\ell-d)\eta&<
 U_0(d)+U_0(x-d+\ell)+U_0(x+\ell)-1+2\varepsilon.\label{eq:triple2norm}
\end{align}
The omitted terms in \eqref{eq:seednorm} are the favorable
$R_0/(2n^2)+1/(4n)$ on the left. The right-hand corrections in
\eqref{eq:triple1norm} and \eqref{eq:triple2norm} are respectively
$(2x+\ell-2)/(2n)<0$ and $(x+\ell-3/2)/n<0$.

If $r^2/2\ge\varepsilon$, \eqref{eq:smallV} is impossible.
Otherwise, when $P<D_a$, geometry gives
\begin{equation}\label{eq:winggeometry}
 G\subset\Uset,\qquad |\Cset\cap\Uset|\ge|G|,\qquad
 |\Cset\cap I|\ge(|G|-P)_+.
\end{equation}
To see this, let $a_L$ be the end of $L$. A left protrusion has
$a_L<t$; at least $\min(D_a,t-a_L)$ wing positions lie outside $L$.
Thus $P<D_a$ forces $a_L>s$ and $t-a_L\le P$.
By \eqref{eq:coremass} the other end of $G$ is before the other end of
$I$, proving \eqref{eq:winggeometry}. It also covers a gap wholly in
the wing. Additional exterior pieces of $\Cset$ only increase the
intersection counts. No such conclusion is used when $P=D_a$.

\subsection{The coupled retained-wing inequality}
Choose rational $d_0$ with $0<d_0\le b(\alpha)-v\le d$ and write
\[
 X=\beta,\qquad \mathfrak b=b(X)^2,\qquad D_0=\mathfrak b-d_0^2,
 \qquad r_X=1-X-\ell.
\]
Replacing $x$ by $X$ weakens the necessary inequalities.
In \eqref{eq:seednorm}, $r^2/2-b(x)^2$ decreases with $x$;
the right side of \eqref{eq:triple1norm} increases with $x$;
the derivative of left minus right in \eqref{eq:triple2norm} is
$\eta-2(1-x-\ell)-d\le0$, since $\eta\le d$.
The derivative of the resulting gap lower bound is
$-z-2(1-x)<0$. The analogous derivatives for
\eqref{eq:full1}--\eqref{eq:full2} are $-z-2(1-x)$ and $-z-(1-x)$.

If $r_X^2\ge2\varepsilon$, no protrusion on either side is possible;
use \eqref{eq:full1}--\eqref{eq:full2} at $x=X,d=d_0$, as their left
sides increase with $d$. Otherwise \eqref{eq:epscond} implies
\begin{equation}\label{eq:elllarge}
 \ell>1-X-\sqrt{2\varepsilon}\ge X>b(X)>d.
\end{equation}
Here $X^2-b(X)^2=2(X-1/2)^2>0$. In particular a protrusion forces
$j>k$; shorter low intervals occur only in the complete-containment
case, which needs no restriction on $j$.

For $A>0$ put $\rho(A,N)=0$ when $N\le0$, and otherwise put
$\rho(A,N)=(\sqrt{A^2+4N}-A)/2$. Define
\begin{align}
 t_S&=\rho(\ell-d_0,\varepsilon-r_X^2/2+D_0),\notag\\
 t_1&=\rho(\ell,U_0(X)+U_0(X+\ell)-3/4+\varepsilon),\label{eq:roots}\\
 t_2&=\rho(X+\ell-d_0,
 U_0(d_0)+U_0(X-d_0+\ell)+U_0(X+\ell)-1+2\varepsilon).\notag
\end{align}
Let $t_* =\min(t_S,t_1,t_2)$.
\begin{lemma}[Coupled retained wing]\label{lem:coupled}
If $t_*<d_0$, then $q=d-\eta\ge d_0-t_*>0$.
\end{lemma}
\begin{proof}
Inequality \eqref{eq:triple1norm} gives $\eta<t_1$ (or is impossible
if its right side is nonpositive). For \eqref{eq:seednorm}, substitute
$\eta=d-q$ to obtain the necessary inequality
\[
 F_S(d,q)=d^2+\ell d-(d+\ell)q+q^2+r_X^2/2-\mathfrak b
 <\varepsilon.
\]
On $q\le d_0$, $d\ge d_0$, this increases with $d$; at $d=d_0$
it decreases with $q$ because $2q-d_0-\ell<0$ by
\eqref{eq:elllarge}. If $t_S<d_0$ and $q<d_0-t_S$, it is at least
$\varepsilon$, a contradiction. This also covers $t_S=0$, since
then $F_S(d_0,d_0)\ge\varepsilon$.

The identity $U_0(d)+U_0(X+\ell-d)=U_0(X+\ell)+(X+\ell)d-d^2$
turns \eqref{eq:triple2norm} into
\[
 F_2(d,q)=d^2-(X+\ell+d)q+q^2<2U_0(X+\ell)-1+2\varepsilon.
\]
Its derivative in $d$ is $2d-q>0$; at $d=d_0$ its derivative in $q$
is $2q-X-\ell-d_0<0$. The same argument gives $q\ge d_0-t_2$,
including the zero-root case. Use whichever root attains $t_*$.
This bounds the retained wing $q$, not $\eta$ at a substituted
diameter.
\end{proof}

When $q_0=d_0-t_*>0$, the wing is nonsaturated and
\eqref{eq:winggeometry} applies. Substitution into \eqref{eq:quad1}
gives, for $z\ge\eta$, exactly
$z^2+(1-x-h+q)z+h^2/2-b(x)^2+q^2$.
For $z<\eta$, the true expression exceeds this by
$(d+q)(\eta-z)\ge0$. It increases in $q$, with derivative $z+2q$.
Thus in both cases
\begin{equation}\label{eq:coupledgap}
 z^2+(1-X-h+q_0)z+h^2/2-\mathfrak b+q_0^2<2\varepsilon.
\end{equation}
If $G\subset I$, \eqref{eq:full1} at $X$ and the actual $d$ implies
the same inequality since $q_0\le d_0\le d$ and its left side
increases with that parameter. Right protrusions follow by reflection.
If $t_*\ge d_0$, only the crude bound $z\le h$ is used.

\subsection{Counting at the lower threshold}
Let $M_j$ count length-$j$ intervals with sum $>Q$. Put
$r_\varepsilon^2=1/2+2\varepsilon$ and
\[
 g_\varepsilon(x)=\begin{cases}0,&x<1-r_\varepsilon,\\
 \min\{1-x,\sqrt{r_\varepsilon^2-(1-x)^2}\},&x\ge1-r_\varepsilon.
 \end{cases}
\]
For overlapping or adjacent extreme intervals, \eqref{eq:pair}
bounds their separation by $ng_\varepsilon(j/n)$; the remaining
finite correction is $j-n\le0$. Using one interval twice excludes
$j/n<1-r_\varepsilon$. Otherwise \eqref{eq:complementmass} applies
and the separation is $n(\ell+z)$. For $y=1/2-j/n\in[0,\sqrt\varepsilon]$,
\begin{equation}\label{eq:rowcount}
 M_j\le ng_\varepsilon(j/n)+1+nF(y),\qquad
 F(y)=(K_{\rm span}(y)-b_\varepsilon(y))_+,
\end{equation}
where $b_\varepsilon(y)=\sqrt{1/4+2\varepsilon-y-y^2}$.
An upper bound $K_{\rm span}$ is the minimum of applicable caps:
\begin{align*}
 &1/2+y;\\
 &\ell+\rho(1-X+d_0-2y,2\varepsilon-2y^2+D_0),\quad
  \ell+\rho(1-X-2y,2\varepsilon-2y^2+D_0/2)
     &&(r_X^2\ge2\varepsilon);\\
 &\ell+\rho(1-X+d_0-2y-t_*,
   2\varepsilon-2y^2+D_0+2d_0t_*-t_*^2)
     &&(r_X^2<2\varepsilon, t_*<d_0).
\end{align*}
Outside this $y$ range no $F$ is needed. If $g_\varepsilon$ uses the
total-start cap then $b_\varepsilon\ge1/2+y\ge K_{\rm span}$, so
$F=0$ and \eqref{eq:rowcount} remains valid.
As before, $g_\varepsilon$ has variation less than $2$ and integral
$\pi/16+\pi\varepsilon/4$. Suppose an upper rectangle sum $B_i$
for $F$ uses grid $M_{\rm grid}=100000$. Each grid cell contains at
most $n/M_{\rm grid}+1$ possible $y$ values. Counting low values by
$Q$ and high values by intervals therefore gives
\begin{equation}\label{eq:lowerthreshold}
 |\Sset(a)|\le\{B_0-(1-\pi/4)\varepsilon+B_i\}n^2
       +\{7/2+M_{\rm grid}B_i\}n.
\end{equation}
The linear terms are $n/2$ from $Q$, at most $2n$ from the Riemann
error, $n$ from the per-length allowance, and $nM_{\rm grid}B_i$
from grid occupancy.

\subsection{Certificate and exact verification}\label{sec:certificate}
The electronic certificate contains $892$ disjoint bins. A row
$[a_i,b_i,v_i',e_i']$ means
\[
 [\alpha_i,\beta_i)=[a_i/10000,b_i/10000),\quad
 v_i=v_i'/10^6,\quad\varepsilon_i=e_i'/4000.
\]
They range from $2929/10000$ to $3821/10000$. The verifier checks their
order and disjointness. We describe all directed-rounding conventions
so that the finite certificate has a precise mathematical meaning.

At scale $R_*=10^{12}$, round $b(\alpha_i)$ down and set the rational
$d_0=b(\alpha_i)_{\rm down}-v_i$. Round $b(\beta_i)$ up for
\eqref{eq:epscond}, and round $D_0$ up. Subsequent formulas use this
particular rational $d_0$; upward rounding of $D_0$ weakens the bounds.
The positive exterior lower bound also has a direct rational check:
if $x=R_*\beta_i$ and $b_+$ is the integer upward root, verify
$R_*-x-b_+>0$ before squaring. The preserved certificate program
does not assert that sign separately; the supplemental publication
check does, and the positivity follows analytically above.

For a cell $[y_0,y_1]=[r/M_{\rm grid},(r+1)/M_{\rm grid}]$, start
with cap $1/2+y_1$. If $(1/2+y_0-X)^2\ge2\varepsilon$, use the two
complete-containment roots. Otherwise bound the three roots in
\eqref{eq:roots} upward, using denominator coefficients at $y_1$ and
numerators at $y_0$. The numerators decrease with $y$: $r_X$ increases,
while $\ell$ and the arguments of the increasing function $U_0$ decrease.
The denominator coefficients decrease as well. If the minimum upward
root $t_{\rm up}<d_0$, use the partial cap; otherwise keep the crude
cap. A cell crossing the complete-containment boundary may still use
the partial cap, since on its complete part the retained parameter is
$d\ge d_0\ge d_0-t_{\rm up}$.

In the gap root use $h=2y_1$ in the coefficient and
$2\varepsilon-2y_0^2$ in the right side. The term
$2d_0t-t^2$ increases for $0\le t<d_0$, so use $t_{\rm up}$.
Outside the root use $\ell\le1/2-y_0$. The preserved program further
replaces the gap root by its maximum with $t_{\rm up}$, a conservative
extra allowance. Subtract $b_\varepsilon(y_1)$ rounded down and take
the positive part, obtaining an upper bound $U_{ir}$ on the cell.
The final cell may extend past $\sqrt\varepsilon$; the bound is needed
only on the intersection. Set $B_i=M_{\rm grid}^{-1}\sum_r U_{ir}$.

All substantive arithmetic is integer arithmetic. For example the
scaled function is $U(v)=v-\lfloor v^2/(2R_*)\rfloor$, an upward
bound; every square root is explicitly rounded up or down, every
positive-denominator division is directed, negative root numerators
are replaced by zero, and every root coefficient is asserted positive.
Thus no floating-point value decides a certificate inequality.
The resulting exact checks are
\begin{align}
 W_{\rm def}:=\sum_i(\beta_i-\alpha_i)v_i&=0.002798309,
 &\sum_i v_i&=27.98309,\label{eq:Wdef}\\
 B_0-W_{\rm def}&\le C_*,
 &C_*&=4435834/10^7,\label{eq:defcheck}\\
 B_0-(1-\pi/4)\varepsilon_i+B_i&\le C_*&&\text{for all }i,\label{eq:rowcheck}\\
 9/4+\sum_i v_i&\le1100,
 &7/2+M_{\rm grid}B_i&\le1100\quad\text{for all }i.\label{eq:lincheck}
\end{align}
There are $17\,430\,410$ rectangles: $8\,938\,271$ partial,
$1\,221\,306$ crude-cap and $7\,270\,833$ complete cells.
For orientation, the largest $B_i$ is approximately
$0.01058603280678933$, the largest coefficient in \eqref{eq:rowcheck}
is approximately $0.4435833245825063$, and the deficit coefficient is
approximately $0.44355123184936207$. These decimals are descriptive;
\eqref{eq:defcheck}--\eqref{eq:lincheck} use integer cross multiplication.

The verifier proves $3.141592653589793<\pi<3.141592653589794$ by
Machin's identity $\pi=16\arctan(1/5)-4\arctan(1/239)$, using $40$
and $12$ alternating-series terms and directed first-omitted-term
bounds. Tangent addition gives $\tan(2\arctan(1/5))=5/12$,
$\tan(4\arctan(1/5))=120/119$, and
$\tan(4\arctan(1/5)-\arctan(1/239))=1$; the last angle is in
$(0,\pi/2)$, proving the identity.

\subsection{Dichotomy and the linear term}
If no bin supplies $e_k\le v_i$, each bin contains at least
$n(\beta_i-\alpha_i)-1$ integer lengths. Disjointness and $e_k\ge0$
give $\sum_ke_k\ge nW_{\rm def}-\sum_iv_i$, including bins whose
lower count is negative. Equations \eqref{eq:baseline},
\eqref{eq:defcheck}, and \eqref{eq:lincheck} prove
$|\Sset(a)|\le C_*n^2+1100n$. Otherwise an anchor pair exists,
and \eqref{eq:lowerthreshold}, \eqref{eq:rowcheck}, and
\eqref{eq:lincheck} prove the same bound. This holds for every $a$
and every positive integer $n$.
Finally $166\cdot66\,265\,061=11\,000\,000\,126\ge11\,000\,000\,000$,
so $1100/n\le0.0000166$ for $n\ge66\,265\,061$.
This proves Theorem~\ref{thm:upper}.

\section{The lower construction and its base}\label{sec:construction}
Let $p=2k+1$ be prime, $N=p-1$, $T=(9p+3)/2$, and $g$ the least
primitive root modulo $p$. For clarity, a finite subgroup of the
multiplicative group of a field is cyclic: take, for each prime
dividing its exponent, an element realizing the maximal prime-power
part, and multiply the resulting coprime-order components. An element
of order the exponent results. The polynomial root bound applied to
$X^e-1$, and Lagrange's theorem, force this order to be the group size.
Thus $g$ exists. Write $[x]_p$ for the least nonnegative residue and set
\begin{equation}\label{eq:labels}
 x_i=g^i\quad(0\le i<N),\qquad x_N=0.
\end{equation}
For $c\in\Fp^*$ and $t\in\Fp$ define
\begin{equation}\label{eq:construction}
 A_i=[cx_i+t]_p+1,\quad C_i=[cx_i+t+k]_p+1,\quad
 Y_i=p+C_i,\quad Z_i=T-A_i-Y_i.
\end{equation}
The sequence $a(p,c,t)$ concatenates $(A_i,Z_i,Y_i)$ for
$i=0,\ldots,p-1$, including the final zero label.
Both $A$ and $C$ permute $1,\ldots,p$. If $A_i\le k+1$, then
$C_i=A_i+k$ and $Z_i=6k+5-2A_i$; otherwise
$C_i=A_i-k-1$ and $Z_i=8k+6-2A_i$. These ranges give respectively
the odd and even integers in $2p+1,\ldots,3p$. Thus $a(p,c,t)$ is a
permutation of $1,\ldots,3p$, and every triple sums to $T$.

Its prefix sums are
\[
 U_j=jT\ (0\le j\le p),\quad V_i=iT+A_i,\quad
 W_i=(i+1)T-p-C_i\ (0\le i<p).
\]
The positive chronological differences having an endpoint $U$ form
a base set $\Bset_0$ with
\begin{equation}\label{eq:basecount}
 |\Bset_0|=2p^2+3p.
\end{equation}
Here is a direct count. Near $hT$, the central offsets are $+A_j$
for $j\ge h$, $0\le h<p$, and $-A_i$ for $i\le p-h$, $1\le h\le p$;
indices are always in $[0,p-1]$. They give $p(p+1)$ values.
In each outer interval $(hT+p,(h+1)T-p)$, $0\le h<p$, the values are
\begin{equation}\label{eq:basetails}
 hT+p+C_i\quad(i\le p-1-h),\qquad
 (h+1)T-p-C_j\quad(j\ge h).
\end{equation}
These give another $p(p+1)$ values. Their two ranges are disjoint
because $T>4p$. The central bands $[hT-p,hT+p]$ and the outer
intervals are disjoint. Finally the $p$ positive multiples of $T$
are distinct from all these values, proving \eqref{eq:basecount}.

Use the central candidates
\begin{equation}\label{eq:centralcandidates}
 dT+A_{i+d}-A_i,\qquad dT-C_{i+d}+C_i
 \quad(1\le d<N, 0\le i<N-d).
\end{equation}
Each class is injective at fixed $d$, because its residue coefficient
in $g^i$ is $\pm c(g^d-1)\ne0$. Distinct central bands are disjoint.
Use also the cross candidates
\begin{equation}\label{eq:crosscandidates}
 R_{d,i}=(d+1)T-p-C_{i+d}-A_i,\qquad
 Q_{d,i}=(d-1)T+p+C_i+A_{i+d},
\end{equation}
for $2\le d<N$, $d\ne k$, $0\le i<N-d$. The coefficients
$\mp c(1+g^d)$ are nonzero precisely when $d\ne k$, so each class
is injective in its outer interval. For $2\le d<N$ there are
$(N-1)(N-2)$ cross candidates; deleting $d=k$ removes exactly $N$.
Their number is $p^2-6p+7\ge p^2-6p$.
The omitted lag $1$ costs only a linear term and makes later domains
uniform. Central and cross regions are disjoint.

\section{Fourier completion and the new central values}\label{sec:fourier}
\subsection{Elementary Fourier norms}
On a finite abelian group use normalized Fourier coefficients and
inversion $f(x)=\sum_\chi\widehat f(\chi)\chi(x)$, and define
$\normA f=\sum_\chi|\widehat f(\chi)|$. Convolution and the triangle
inequality show that $\normA{fg}\le\normA f\normA g$.
Affine pullback cannot increase this norm, since it only combines
coefficients and attaches unit phases. Put $L=\log p+3$.

On $\mathbb Z/q\mathbb Z$, $q\le p$, any cyclic interval has norm
at most $L$. Its zero coefficient is at most $1$, and geometric
summation with $\sin(\pi r/q)\ge2r/q$, $0\le r\le q/2$, bounds the
coefficient at cyclic distance $r>0$ by $1/(2r)$. Summation gives
$1+H_{\lfloor q/2\rfloor}\le2+\log q\le L$.
For representatives $x,y\in[0,q)$, both $\ind_{x<y}$ and
$\ind_{x+y<q}$ have norm at most $2L$. Indeed, put $s(x)=[x]_q/q$.
Then
\[
 \ind_{x+y<q}=1-s(x)-s(y)+s(x+y),\qquad
 \ind_{x<y}=s(y)+s(q-1-x)-s(q-1-x+y).
\]
Every nonconstant coefficient of $s$ has modulus
$1/(q|1-e^{2\pi ir/q}|)\le1/(4|r|_q)$, so its nonconstant mass
is at most $H_{\lfloor q/2\rfloor}/2$. The displayed identities
have constant term at most $1$ and nonconstant mass at most
$3H_{\lfloor q/2\rfloor}/2$, proving the bound.

An ordered triangle in an ordinary interval $J$ has indicator
$\ind_J(u)\ind_J(v)\ind_{u<v}$ and hence norm at most $2L^3$.
This remains true after independent translations of $u,v$.
For $1\le a,b\le p$, the triangle $a+b\le p$ is a reflected
comparison, with norm at most $2L$. For
$a+b\ge(3p+3)/2$, restrict $a\ge(p+3)/2$ and use
$u=a-(p+3)/2$, $v=p-b$ and the comparison $v\le u$;
the bound $2L^3$ is more than sufficient. We use the common bound
$8L^3$ for all the two-variable indicators below.

\subsection{A graph-completion lemma}
Write $e_q(v)=\exp(2\pi iv/q)$. For an indicator $R$ on $\Fp^2$,
write $|R|=\sum_{X,Y}R(X,Y)$.
\begin{lemma}\label{lem:graph}
Let $\alpha,\beta\in\Fp^*$, $s\in\Fp$, and
$F(X,Y,Z)=R(X,Y)\ind_{Z=\alpha X+\beta Y+s}$, with
$\normA R\le K_R$. For $D\subset(\mathbb Z/N\mathbb Z)^3$,
\begin{equation}\label{eq:graph}
 \left|\frac1N\sum_{c\ne0}\sum_{(d,i,j)\in D}
 F(cg^i,cg^{i+d},cg^j)-\frac{|D|\,|R|}{p^3}\right|
 \le K_R(p^{3/2}+3p)\normA{\ind_D}.
\end{equation}
\end{lemma}
\begin{proof}
The Fourier coefficient at $(u,v,w)$ is
$p^{-1}e_p(-ws)\widehat R(u+\alpha w,v+\beta w)$.
Its total mass is $\normA R$, and its mass on each of $u=0$, $v=0$,
$w=0$ is $\normA R/p$: the relevant maps to the two arguments of
$\widehat R$ are bijections, using $\alpha\beta\ne0$.
Let $\chi_a(g^r)=e_N(ar)$ and
$G(\chi,u)=\sum_{x\ne0}\chi(x)e_p(ux)$. Completing an index
character gives exactly
\begin{align}\label{eq:gausscompletion}
 &\frac1N\sum_{c\ne0}\sum_{d,i,j\bmod N}
 e_N(ad+bi+ej)e_p(c(ug^i+vg^{i+d}+wg^j))\notag\\
 &\hspace{20mm}=\ind_{b+e=0}
 G(\chi_{b-a},u)G(\chi_a,v)G(\chi_e,w).
\end{align}
Substitute $X=cg^i$, $Y=cg^{i+d}$, $Z=cg^j$; the residual character
in $c$ is $\chi_{-b-e}(c)$, proving the identity.
At zero additive frequency a Gauss sum is $N$ for the principal
character and zero otherwise. The principal sum at nonzero frequency
is $-1$. For nonprincipal $\chi$, rescaling to $u=1$ gives
\[
 |G(\chi,1)|^2=\sum_{z\ne0}\chi(z)
 \sum_{y\ne0}e_p((z-1)y)
 =N-\sum_{z\ne0,1}\chi(z)=p.
\]
Thus \eqref{eq:gausscompletion} is at most $p^{3/2}$ in modulus when
$u,v,w$ are nonzero. If exactly one is zero it is at most $Np\le p^2$.
If exactly two are zero, the condition $b+e=0$ forces all three
characters to be principal and the modulus is at most $N^2\le p^2$.
If all three are zero, only $a=b=e=0$ contributes, with value $N^3$.
Fourier inversion gives the principal term $|D||R|/p^3$.
The all-nonzero frequencies cost at most
$K_Rp^{3/2}\normA{\ind_D}$; the other frequencies cost at most
$3K_Rp\normA{\ind_D}$, using the three coordinate-plane masses.
\end{proof}

Every three-label index domain used below satisfies
\begin{equation}\label{eq:domainnorm}
 \normA{\ind_D}\le16L^5.
\end{equation}
Here are explicit factorizations. For representatives modulo $N$ put
$U(a,b)=\ind_{[a]_N+[b]_N<N}$ and
$T_<(a,b)=\ind_{[a]_N<[b]_N}$. Positive missing labels use
$T_<(j,d)U(i,d)$; negative missing labels use
\[
 \ind_{[3,N-1]}(d)U(i,d)T_<(N-1-j,d-2).
\]
The latter says $j\ge N-d+2$ and has norm at most $4L^3$;
the former has norm at most $4L^2$.
The central-overlap domain has indicator
\begin{equation}\label{eq:overlapdomain}
 T_<(j,d)U(i,d)\ind_{[0,k-1]}(d)
 \bigl(\ind_{[0,k-1]}(i+d)+\ind_{[k,N-1]}(i)\bigr).
\end{equation}
The no-wrap factor makes the two summands disjoint. Its norm is at
most $8L^4$; reflecting $j$ preserves the bound.
Cross-base domains have $d\in[2,N-1]$, $U(i,d)$, and one of
$j\ge d$, $U(d-1,j)$, $U(d-2,j)$, $j\ge d-1$.
Comparison complements have norm at most $1+2L\le3L$, so these
domains have norm at most $6L^3$. These prove \eqref{eq:domainnorm}.
For $p\ge2^{16}$, $p^{3/2}+3p\le2p^{3/2}$. With $K_R\le8L^3$,
each such count has error at most
\begin{equation}\label{eq:Eerror}
 E_0=256p^{3/2}L^8.
\end{equation}

\subsection{Separate central novelty, uniformly in translation}
At lag $d$, positive offset $A_j$ is missing from $\Bset_0$ exactly
when $j<d$; negative offset $-A_j$ is missing exactly when $j>p-d$.
For the negative sign use only $j\in[N-d+2,N-1]$, an empty set when
$d<3$. Zero-label missing values are unused in this lower count.
For either class in \eqref{eq:centralcandidates}, equality to a signed
missing label is a graph as in Lemma~\ref{lem:graph}, with coefficients
$\alpha,\beta\in\{1,-1\}$. Its two-variable indicator specifies the
sign of a difference in $1,\ldots,p$, has mass $p(p-1)/2$, and obeys
the norm bound above, for every $t$. For example
$A(Y)-A(X)=A(Z)$ gives $Z=Y-X-t-1$ and $A(Y)>A(X)$.
The other class uses $C(X)-C(Y)$.

The domain sizes are
\[
 M_+=\sum_{d=1}^{N-1}(N-d)d=\frac{N(N^2-1)}6,\quad
 M_-=\sum_{d=3}^{N-1}(N-d)(d-2)=\frac{(N-2)((N-2)^2-1)}6.
\]
Their sum is $p^3/3-2p^2+14p/3-4\ge p^3/3-2p^2$.
Thus each class has expected new count, averaging $c\ne0$ with $t$
fixed, at least
\[
 \frac{p-1}{2p^2}(p^3/3-2p^2)-2E_0
 \ge p^2/6-3p/2-2E_0.
\]
There is a unique signed missing label for each value, and class
injectivity prevents repetitions. Adding the two classes yields
\begin{equation}\label{eq:separatecentral}
 \E_c(\text{separate central new counts})\ge p^2/3-3p-4E_0.
\end{equation}

\subsection{Restricted translations and central overlap}
Set
\begin{equation}\label{eq:translationinterval}
 m=\lfloor p^{3/4}\rfloor,\qquad I_p=\{k-m,\ldots,k+m\},\qquad M=2m+1.
\end{equation}
For $p\ge2^{16}$ this is an ordinary interval in $[0,p-1]$ and
$2m\le k-1$. Indeed $p^{1/4}\ge16$ implies $2p^{3/4}\le p/8$,
and $p/8\le(p-3)/2$ for $p\ge4$. From now on averages are over
$c\in\Fp^*$ and $t\in I_p$.

Equality between the central classes at the same lag forces the
second start label to be $-g^i$. Thus $d<k$ and the possible starts are
\begin{equation}\label{eq:centralstarts}
 i\in[0,k-d-1]\ \cup\ [k,N-d-1].
\end{equation}
Write $X=cg^i$, $Y=cg^{i+d}$, $t=k+e$, $u=[X+t]_p$,
$v=[Y+t]_p$, and $b=[2t+k]_p$. Since $C(-X)=[b-u]_p+1$,
the lifts agree exactly when $u,v$ lie in the same one of $[0,b]$
and $[b+1,p-1]$; the common lift is positive exactly when $u<v$.
The condition $2|e|\le k-1$ gives $b=k-1+2e$.
The positive and negative common-lift indicators $R_+,R_-$ on the
whole field square have equal masses
\begin{align}\label{eq:liftmass}
 |R_\pm|&=\binom{k+2e}{2}+\binom{p-k-2e}{2}\notag\\
 &=p^2/4-p/2+4e^2-2e+1/4
 \le p^2/4+4m^2+2m+1.
\end{align}
Each is the sum of two same-side ordered triangles; consequently
$\normA{R_\pm}\le4L^3$.

A positive common lift equal to a missing $A_j$ gives
$Z=Y-X-t-1$; a negative one gives $Z=X-Y-t-1$.
For an upper bound enlarge the nonzero negative missing-label set to
$j\in[N-d,N-1]$. Both signs now use \eqref{eq:overlapdomain}, or its
reflection, with domain size
\[
 \sum_{d=1}^{k-1}d(N-2d)=(k^3-k)/3\le p^3/24.
\]
The zero label can occur only for the negative missing sign and
contributes at most $p$ common values over all lags.
Lemma~\ref{lem:graph} and \eqref{eq:liftmass} give
\begin{align}\label{eq:centraloverlap}
 \E(\text{central overlap outside }\Bset_0)
 &\le p^2/48+(4m^2+2m+1)/12+p+2E_0\notag\\
 &\le p^2/48+p^{3/2}+p+2E_0.
\end{align}
Here $m\ge1$ implies $4m^2+2m+1\le7m^2\le12p^{3/2}$.
Subtracting \eqref{eq:centraloverlap} from \eqref{eq:separatecentral}
proves
\begin{equation}\label{eq:centralunion}
 \E(\text{new central union})\ge5p^2/16-4p-p^{3/2}-6E_0.
\end{equation}

\section{Cross terms and restoration}\label{sec:cross}
\subsection{Separate cross novelty}
Put $H=T-2p=(5p+3)/2$. A collision of either cross class with the
nearer base tail has form $a+b=C_j$, and one with the opposite tail
has form $a+b+C_j=H$, where $(a,b)$ is $(A_i,C_{i+d})$ or
$(C_i,A_{i+d})$. Each congruence is a three-label graph with both
coefficients nonzero. The ordinary triangle has mass $p(p-1)/2$.
The opposite triangle $a+b\ge H-p$ has mass
$k(k+1)/2=(p^2-1)/8$: reflect from $(p,p)$ to nonnegative coordinates
whose sum is at most $k-1$. Both obey the established norm bound.

For $R$ the nonzero ordinary base indices are $j\ge d$ and the
opposite indices $j\le N-d$. For $Q$ they are respectively
$j\le N+1-d$ and $j\ge d-1$, always clipped to $[0,N-1]$.
These are the four domain forms above. Include $d=k$ for upper bounds.
For each of the ordinary and opposite totals, the two domain sizes sum
to at most
\begin{equation}\label{eq:crossdomainsize}
 \sum_{d=1}^{N-1}(N-d)(2(N-d)+2)
 =\frac{2(N-1)N(N+1)}3\le\frac23p^3.
\end{equation}
The zero label costs at most $4p$ collisions, by injectivity (one per
lag, class and tail type). The ordinary loss is at most $p^2/3+2E_0$,
and the opposite loss at most $p^2/12+2E_0$, by Lemma~\ref{lem:graph}.
Subtracting these and the zero-label allowance from at least
$p^2-6p$ candidates proves, uniformly in $t$,
\begin{equation}\label{eq:separatecross}
 \E(\text{separate cross new counts})\ge7p^2/12-10p-4E_0.
\end{equation}

\subsection{Four and five labels under restricted translation}
Put $q=4p$, $L_0=\log p+2$, and, for $r\in\mathbb Z/q\mathbb Z$,
\[
 K(r)=p^{-1}\sum_{x=1}^p e_q(rx),\quad
 B(r)=\min(1,2/|r|_q)\ (r\ne0),\quad B(0)=1,
 \quad B_1=\sum_rB(r),\quad B_2=\sum_rB(r)^2.
\]
Geometric summation gives $|K(r)|\le B(r)$, while
\begin{equation}\label{eq:Bbounds}
 B_1\le1+4H_{2p}\le5+4\log(2p)\le4L_0,
 \qquad B_2<15.
\end{equation}
The last inequality follows from
$B_2\le1+8\sum_{j\ge1}j^{-2}$ and
$\sum_{j\ge1}j^{-2}<1+1/4+\int_2^\infty x^{-2}\,dx=7/4$.
For representatives $1\le x_j\le p$ define
\[
 F_4(x)=\ind_{x_1+x_2+x_3+x_4=H},\qquad
 F_5(x)=F_4(x_1,x_2,x_3,x_4)\ind_{x_3+x_4=x_5}.
\]
Congruences modulo $q$ impose exactly these integer equalities,
because every possible nonzero discrepancy has magnitude less than
$q$. Expanding them gives
\begin{align}
 \widehat F_4(\xi)&=q^{-1}\sum_a e_q(-aH)
 \prod_{j=1}^4K(a-4\xi_j),\label{eq:F4hat}\\
 \widehat F_5(\xi)&=q^{-2}\sum_{a,b}e_q(-aH)
 K(a-4\xi_1)K(a-4\xi_2)\notag\\
 &\hspace{14mm}\cdot K(a+b-4\xi_3)K(a+b-4\xi_4)K(-b-4\xi_5).
 \label{eq:F5hat}
\end{align}
For every $s\in\Fp$,
\begin{equation}\label{eq:planes}
 \sum_{\sum\xi=s}|\widehat F_4(\xi)|\le B_1^4/p,
 \qquad
 \sum_{\sum\xi=s}|\widehat F_5(\xi)|\le B_1^5/q.
\end{equation}
For \eqref{eq:F4hat} put $r_j=a-4\xi_j$. A fixed tuple $r$ obeys
$\sum r_j=4a-4s\pmod q$ and allows at most four $a$; for each,
the $\xi_j$ are unique if they exist. For \eqref{eq:F5hat} use its
five arguments $r_j$; then $\sum r_j=4a+b-4s$, so there are at most
$q$ choices of $a$, with $b,\xi$ determined. This proves the bounds.
On the exceptional subspaces, still for every plane $\sum\xi=s$,
\begin{align}
 \sum_{\sum\xi=s,\ \xi_5=0}|\widehat F_5|&\le4B_1^5/q^2,
 \notag\\
 \sum_{\sum\xi=s,\ \xi_1=\xi_2=0}|\widehat F_5|&\le B_2B_1^3/q^2,
 \label{eq:exceptionalplanes}\\
 \sum_{\sum\xi=s,\ \xi_3=\xi_4=0}|\widehat F_5|&\le3B_2B_1^3/q^2.
 \notag
\end{align}
In the first case $b=-r_5$ and $4a=\sum r_j+r_5+4s$, with at most
four roots. In the second $r_1=r_2=a$ and
$b=r_3+r_4+r_5-2a+4s$; summing the repeated factor gives $B_2$.
In the third put $w=a+b=r_3=r_4$; then
$3a=r_1+r_2+w+r_5+4s$. There is exactly one root since
$\gcd(3,4p)=1$ for $p>3$. We retain the conservative factor $3$
in \eqref{eq:exceptionalplanes}.

The translation weight $w_I(t)=p/M$ on $I_p$, zero elsewhere, has
mean one and
\begin{equation}\label{eq:weightnorm}
 W_I:=\normA{w_I}\le(p/M)L\le p^{1/4}L.
\end{equation}
Indeed $2\lfloor p^{3/4}\rfloor+1\ge p^{3/4}$.
Translation averaging attaches $\widehat w_I(-\sum\xi)$ to a
Fourier character. Thus multiplying the plane estimates
\eqref{eq:planes}--\eqref{eq:exceptionalplanes} by these absolute
weights and summing over $s$ costs at most the factor $W_I$.
The zero frequency keeps weight one. This accounts for every
nonzero total-frequency plane; no independence or fixed-translation
equidistribution is assumed.

For an interval $J$ of power indices put
$S_J(v)=\sum_{j\in J}e_p(vg^j)$. Interval completion with the Gauss
sums in \eqref{eq:gausscompletion}, and additive orthogonality, gives
\begin{equation}\label{eq:incomplete}
 |S_J(v)|\le\sqrt p L_0\quad(v\ne0),\qquad
 \sum_{v\in\Fp}|S_J(v)|^2=p|J|.
\end{equation}
For the first, the cyclic interval norm is at most
$1+H_{N/2}\le L_0$ and completed Gauss sums have modulus at most
$\sqrt p$. The second uses the distinctness of $g^j$.

\subsection{Total cross overlap}
Only $R$ at lag $d$ and $Q$ at lag $d+1$ share an outer interval.
Their equality is
\begin{equation}\label{eq:crossintersection}
 A_i+C_{i+d}+C_u+A_{u+d+1}=H.
\end{equation}
For an upper bound include every $0\le d\le p-3$,
$0\le i\le p-2-d$, $0\le u\le p-3-d$. The tuple count satisfies
\[
 M_4\le\sum_{d=0}^{p-3}(p-1-d)(p-2-d)
 =p(p-1)(p-2)/3\le p^3/3.
\]
After $c$-averaging the character at $\xi$, apart from a fixed unit
phase, equals
\begin{equation}\label{eq:caverage}
 \frac{p\ind_{(\xi_1+g^d\xi_2)g^i+(\xi_3+g^{d+1}\xi_4)g^u=0}-1}{p-1}.
\end{equation}
For $\xi\ne0$ at most $2p^2$ tuples satisfy the indicated equation.
Unless both coefficients vanish, each $i$ admits at most one $u$,
or none if exactly one coefficient vanishes. Both coefficients can
vanish at at most one lag unless $\xi=0$; that lag contributes at
most $p^2$ tuples and all other lags at most $p^2$.
The summed absolute average is thus at most
$(2p^3+M_4)/(p-1)\le4p^2$ for $p\ge3$.
Equations \eqref{eq:planes}, \eqref{eq:weightnorm}, and
\eqref{eq:Bbounds} bound the nonzero-frequency contribution by
\begin{equation}\label{eq:crossfourerror}
 4W_IpB_1^4\le1024W_IpL_0^4.
\end{equation}
Inclusion-exclusion for four positive summands bounded by $p$ gives
\begin{equation}\label{eq:F4volume}
 |F_4|=\binom{5k+3}{3}-4\binom{3k+2}{3}+6\binom{k+1}{3}
 =\frac{23p^3+15p^2+p+9}{48}.
\end{equation}
Three variables cannot exceed $p$ simultaneously in this count.
Multiplying $|F_4|/p^4$ by $M_4\le p^3/3$ gives at most
$23p^2/144+p$. Hence
\begin{equation}\label{eq:totalcross}
 \E(\text{total cross intersection})
 \le23p^2/144+p+1024W_IpL_0^4.
\end{equation}

\subsection{A five-label estimate}
Let $D$ be a subset of lag indices and let $I_d,U_d,J_d$ be intervals
of power indices, each of size at most $N$. Insert in $F_5$ the
representatives in $1,\ldots,p$ of
\[
 cg^i+t+\kappa_1,\quad cg^{i+d}+t+\kappa_2,\quad
 cg^u+t+\kappa_3,\quad cg^{u+d+1}+t+\kappa_4,\quad
 cg^j+t+\kappa_5,
\]
where the shifts are fixed. Let $M_5=\sum_d|I_d||U_d||J_d|$.
We claim
\begin{equation}\label{eq:fiveestimate}
 \left|\E\sum_{d,i,u,j}F_5-\frac{|F_5|}{p^5}M_5\right|
 \le20000W_Ip^{3/2}L_0^6.
\end{equation}
For a frequency $\xi$ put $a_d=\xi_1+g^d\xi_2$,
$b_d=\xi_3+g^{d+1}\xi_4$, $e=\xi_5$.
If all three are nonzero, \eqref{eq:incomplete} and Cauchy--Schwarz give
\begin{align}\label{eq:threeintervals}
 &\frac1N\sum_{c\ne0}
 |S_{I_d}(ca_d)S_{U_d}(cb_d)S_{J_d}(ce)|\notag\\
 &\qquad\le\sqrt pL_0\frac{p\sqrt{|U_d||J_d|}}N
 \le2p^{3/2}L_0.
\end{align}
In fact the final factor $2$ is unnecessary since
$\sqrt{|U_d||J_d|}\le N$; we retain it throughout.
If exactly two coefficients are nonzero, the two-factor
Cauchy--Schwarz bound with the remaining interval length is at most
$2p^2$. If exactly one is nonzero, average before taking absolute
values: each nonzero field power has character average $-1/N$, so
the modulus is the product of the three lengths divided by $N$,
at most $2p^2$.

Outside the union of the three subspaces in
\eqref{eq:exceptionalplanes}, both $a_d,b_d$ are nonzero polynomials
of degree at most one in $g^d$, and $e\ne0$. At most two lags have
a vanishing coefficient. Thus the summed character contribution is
at most $2p^{5/2}L_0+4p^2$. On that union, for $\xi\ne0$ all
coefficients can vanish at at most one lag: two distinct such lags
would force $\xi=0$. That exceptional lag costs at most $p^3$, and
all other lags at most $2p^2$ each, giving $3p^3$.
Using the plane bounds, the nonzero-frequency error is at most
$W_I$ times
\begin{align*}
 &\tfrac12p^{3/2}L_0B_1^5+\tfrac74pB_1^5
       +\tfrac{3B_2}{4}pB_1^3\\
 &\qquad\le(512+1792+720)p^{3/2}L_0^6
 <20000p^{3/2}L_0^6.
\end{align*}
The zero frequency supplies the stated main term, proving
\eqref{eq:fiveestimate}. Exchanging the first and second coordinate
pairs merely exchanges the exceptional subspaces and obeys the same
bound.

\subsection{Restoring intersections already in the base}
Count $F_5$ by $s=x_3+x_4=x_5$. Necessarily $k+2\le s\le2k+1$;
there are $(s-1)(s-k-1)$ choices for the two pairs. Therefore
\begin{equation}\label{eq:F5volume}
 |F_5|=\sum_{s=k+2}^{2k+1}(s-1)(s-k-1)
 =\frac{k(k+1)(5k+1)}6=\frac{(p^2-1)(5p-3)}{48}.
\end{equation}
Use the retained lags $d\in\{2,\ldots,p-3\}\setminus\{k-1,k\}$,
$I_d=[0,p-2-d]$, and $U_d=[0,p-3-d]$.
At a cross overlap put $s=C_u+A_{u+d+1}$, $v=A_i+C_{i+d}$.
Then $s+v=H$ and the common sum is $dT+p+s=(d+1)T-p-v$.
It is already in $\Bset_0$ if either
\begin{equation}\label{eq:restorebranches}
 s=C_j,\quad 0\le j\le p-2-d,
 \qquad\text{or}\qquad
 v=C_j,\quad d\le j\le p-2.
\end{equation}
These are valid subsets of \eqref{eq:basetails}, using only nonzero
labels. The first branch deliberately drops the additional valid
index $j=p-1-d$, to keep both branch lengths $|J_d|=p-1-d$.
The branches are disjoint because $s,v\le p$ together contradict
$H>2p$. The outer interval fixes $d$; injectivity fixes the $R$ and
$Q$ indices and the covering base label. Thus each branch restores
distinct values, without repetitions between branches.

Each branch has tuple count
\begin{equation}\label{eq:M5bounds}
 M_5=\sum_d(p-1-d)^2(p-2-d),\qquad
 p^4/4-6p^3\le M_5\le p^4.
\end{equation}
For the lower bound begin with
$\sum_{s=1}^p s^3\ge\int_0^p x^3\,dx=p^4/4$.
First use lags $1\le d\le p-3$, corresponding to $2\le s=p-1-d\le p-2$.
Removing the top two cubes costs at most $2p^3$, and replacing
cubes by $s^3-s^2$ costs at most $p^3$ (the $s=1$ term becomes zero).
Removing lags $k-1,k$ costs at most $2p^3$, and removing lag $1$
at most $p^3$. This proves the lower bound; the upper follows from
at most $p$ terms, each at most $p^3$.
Equation \eqref{eq:F5volume} gives
$|F_5|/p^5\ge5/(48p^2)-1/(6p^3)$.
Using the lower bound on $M_5$ for its positive coefficient and the
upper bound for its negative coefficient, each principal term is at least
\[
 \frac5{48p^2}(p^4/4-6p^3)-\frac{p^4}{6p^3}
 =5p^2/192-19p/24.
\]
Applying \eqref{eq:fiveestimate} twice and using $19p/12\le2p$ proves
\begin{equation}\label{eq:restoration}
 \E(\text{restored base intersections})
 \ge5p^2/96-2p-40000W_Ip^{3/2}L_0^6.
\end{equation}
The accounting identity behind this step is
\[
 |(R\cup Q)\setminus\Bset_0|
 =|R\setminus\Bset_0|+|Q\setminus\Bset_0|-|R\cap Q|
  +|R\cap Q\cap\Bset_0|.
\]
The last term is bounded from below only by the two disjoint branches
just proved; the subtraction uses the enlarged domain in
\eqref{eq:totalcross}.

\section{Explicit errors, onsets and padding}\label{sec:onsets}
\subsection{Combining the counts and selecting parameters}
The base, new central union and new cross union are disjoint.
Add \eqref{eq:basecount}, \eqref{eq:centralunion} and
\eqref{eq:separatecross}, subtract \eqref{eq:totalcross}, and restore
\eqref{eq:restoration}. The quadratic coefficient is
\[
 2+5/16+7/12-23/144+5/96=803/288.
\]
The error is at most
\[
 14p+p^{3/2}+10E_0+1024W_IpL_0^4+40000W_Ip^{3/2}L_0^6.
\]
Using $E_0=256p^{3/2}L^8$, $W_I\le p^{1/4}L$, $L_0\le L$ and
$p,L\ge1$, this is at most
\begin{equation}\label{eq:totalerror}
 (14+1+2560+1024+40000)p^{7/4}L^8
 =43599p^{7/4}L^8\le50000p^{7/4}L^8.
\end{equation}
This proves the average assertion and hence \eqref{eq:mainlower}.
Define $a^{(3p)}$ by evaluating every $c=1,\ldots,p-1$ and every
$t\in I_p$, taking a pair maximizing $|\Sset|$, and breaking ties
lexicographically by $(t,c)$. Compute $m$ as the largest integer with
$m^4\le p^3$. All construction steps and comparisons can use integer
arithmetic. This is an unambiguous finite definition.

To justify the exact positivity statement after Theorem~\ref{thm:prime},
write $u=\log p$. The normalized error in units of $p^2$ is
$50000e^{-u/4}(u+3)^8$. Its logarithmic derivative is
$8/(u+3)-1/4$, positive for $u<29$ and negative for $u>29$.
At $u=16\log2$ its value is greater than $50000\cdot3^8/16>803/288$;
also $16\log2<29$. It tends to zero as $u\to\infty$.
This proves existence, uniqueness and the exact criterion
\eqref{eq:positivity}. The approximate exponent $304.220$ is a
numerical description only; it is not used to prove an onset.
For a fully explicit vanishing-error estimate, when $u\ge3$,
$(u+3)^8\le2^8u^8$ and $e^{u/8}\ge(u/8)^8/8!$, whence
\[
 50000e^{-u/4}(u+3)^8
 \le50000\cdot2^8\cdot8^8\cdot8!\,e^{-u/8}.
\]
In particular no unspecified asymptotic error is needed for the
leading coefficient $803/2592$ at prime lengths.

\subsection{Certified prime-length onsets}
Put $x=\log_2p$. Since $\log p\le\log_2p$, the normalized error is
at most $50000(x+3)^8 2^{-x/4}$.
Its logarithmic derivative is $8/(x+3)-(\log2)/4<0$ for $x\ge416$,
using $\log2>2/3$. One elementary proof of the latter is the positive
series $\log2=2\sum_{j\ge0}(1/3)^{2j+1}/(2j+1)$.
The exact integer comparisons
\begin{align}
 50000\cdot419^8\cdot9000000&<31\cdot2^{104},\label{eq:onset416}\\
 50000\cdot451^8\cdot90000000&<67\cdot2^{112}\label{eq:onset448}
\end{align}
and the identities
\[
 \frac{803}{288}-9\frac{309799}{10^6}=\frac{31}{9000000},\qquad
 \frac{803}{288}-9\frac{3097993}{10^7}=\frac{67}{90000000}
\]
prove $|\Sset(a^{(3p)})|\ge0.309799(3p)^2$ for $p\ge2^{416}$,
and the stronger $|\Sset(a^{(3p)})|\ge0.3097993(3p)^2$ for
$p\ge2^{448}$. This proves Theorem~\ref{thm:prime}.
These onsets are sufficient and are not claimed minimal.

\subsection{Padding using the prime-interval theorem}
Let $n\ge2^{4096}$, $b=\lfloor n/3\rfloor$ and
$x=(9999999/10000000)b$. Since $n\ge30$,
$b\ge3n/10$ and $x>n/4\ge2^{4094}>3275$.
Thus $\log x>4094(2/3)>2700$ and
$1/(2\log^2x)<1/14000000$. The rational inequality
\[
 \frac{9999999}{10000000}\frac{14000001}{14000000}<1
\]
shows that \eqref{eq:P} supplies a prime in $(x,b]$.
Choose the largest prime $p\le b$. Then $p>x>2^{4094}$, in particular
$p\ge2^{448}$. Append $3p+1,\ldots,n$ to $a^{(3p)}$; all old
consecutive sums persist. Also
$3b/n\ge1-2/n\ge999999999/1000000000$, so
\[
 \frac{3p}{n}>
 \frac{9999999}{10000000}\frac{999999999}{1000000000}
 >\frac{99999989}{100000000}.
\]
Using \eqref{eq:onset448}, the padded permutation satisfies
\begin{equation}\label{eq:paddingfinal}
 \frac{|\Sset|}{n^2}\ge
 \frac{3097993}{10000000}\left(\frac{99999989}{100000000}\right)^2
 >\frac{309799}{1000000}.
\end{equation}
The last comparison is exact rational arithmetic. This proves the
all-integer statement. Taking lower and upper limits and using
$1100/n\to0$ proves Theorem~\ref{thm:all}.

\section{Computations and reproduction}\label{sec:computations}
The electronic bundle contains the complete upper table, the exact
certificate program and its per-row output, the finite-check program
and output, and all lower-construction programs and recorded outputs.
The README gives exact commands and identifies every file;
\texttt{SHA256SUMS} records their hashes. The supplied programs and
outputs are preserved byte for byte. A separate small publication
check verifies the positive exterior sign mentioned in
Section~\ref{sec:certificate} and the onset arithmetic.
All programs use the Node.js standard library. The upper certificate
uses arbitrary-precision integers, including directed square roots.
Floating decimals in its output are summaries, not certificate inputs.

\subsection{Upper certificate and finite checks}
The $892$-row, $17\,430\,410$-cell certificate is part of the proof;
the following finite permutation checks are supporting tests.
All $46\,233$ permutations for $1\le n\le8$ were enumerated, giving
\[
 (f(1),\ldots,f(8))=(1,3,6,9,13,19,25,32).
\]
For every nontrivial extreme half-high pair with $k\le n/2$, the
checker considers all disjoint equal-length low pairs, including
lengths below $k$, for the general category inequalities and the
three protrusion rearrangements when applicable.
Additional tests use the tent $1,3,5,\ldots$, followed by descending
even entries, its reversal, and $40$ seeded Fisher--Yates permutations
at each $n\in\{16,32,64,128,256,512\}$.
The seed is hexadecimal $17\mathrm d34210$, and the threshold
numerators are $80,100,120,140,160,180,200,220$ over $4000$.
Three $3000$-swap searches at each $n\in\{32,64,128\}$ retain counts
$375,1399,5335$, with $9003$ objective evaluations per length.

In total the preserved checker processes $46\,488$ permutations,
$34\,699$ anchor pairs, $2\,025\,199$ category tests,
$2\,312\,221$ tests of each protrusion rearrangement,
$12\,828$ nonsaturated geometry cases, $10\,999$ saturated cases,
$404\,887$ low rows and $458$ coupled checks.
The coupled trial loop imposes $j\ge k$; this is adequate for its
protrusion branch because \eqref{eq:elllarge} forces $j>k$.
The complete-containment proof itself imposes no such restriction.
Finite tests cannot replace the universal rearrangement argument.
For this bundle the exact certificate and this checker were replayed;
their outputs agree with the preserved outputs after removing timing
fields.

\subsection{Lower-construction checks}
Direct bitset enumeration covers all $(c,t)$ for
$p=5,7,11,17,31,61,101$, including means and extrema by translation.
A separate rotating-window multiset counter covers every $(c,t)$ for
$p=5,7,11,17,31,101,251,503,1009$, a total of $1\,343\,802$ pairs.
Counts through $p=31$ are cross-checked against the direct counter;
the two independent scans also agree on the maximum and mean at
$p=101$. The maximum counts for $p=251,503,1009$ are
$176\,678$, $707\,901$, and $2\,844\,224$ respectively.
The largest exhaustively scanned prime is $1009$.

The exact audit examines all $3190$ pairs for
$p=5,7,11,13,17,19,23,29,31$. It checks the permutation, base
cardinality, direct bitset count against a separate set enumeration,
class injectivity, missing-label characterization, membership of
candidates among interval sums, and $64\,199$ restored overlap
values with disjoint restoration branches. It also checks $82\,115$
central-lift equivalences, $69\,492$ domain masks, $27$ domain-size
identities and $173$ volume identities, together with
\eqref{eq:onset416}--\eqref{eq:paddingfinal} in arbitrary-precision
integer arithmetic. This audit was replayed for the bundle and agrees
with its preserved output apart from timing.

The recorded samples at $p=251,503,1009,2003,3331$ use multipliers
$1,2,g,p-1$ and translations
$0,\lfloor p/8\rfloor,\lfloor p/4\rfloor,\lfloor3p/8\rfloor,
(p-1)/2,\lfloor5p/8\rfloor,\lfloor3p/4\rfloor,\lfloor7p/8\rfloor$,
for $160$ counts. Padded examples give $30\,981\,553$ sums at
$n=10000$, using $(p,c,t)=(3331,3,1665)$, and $69\,813\,779$ at
$n=15010$, using $(5003,2,2501)$. These are counts for particular
permutations, not the extremal values $f(n)$.
These finite computations test combinatorial identities. They cannot
test the content of \eqref{eq:mainlower} at its enormous onset, and
finite unrestricted maxima above $0.31n^2$ do not prove an asymptotic
gain. The theorem averages over the narrower set $I_p$.

The rotating counter is exact for the tested range. Multiplication
$c=g^r$ rotates the $N$ nonzero triple blocks, while the zero block
stays last. It maintains the multiset of intervals inside that core,
inside the zero block, and joining the core to the zero block.
A one-block rotation removes the core intervals starting in the
three departing positions and adds those ending in the three new
positions, then replaces the joining intervals. Other sums persist.
Changes of multiplicity across zero update the distinct-value count.
The $N$ rotations visit each nonzero multiplier once. This reduces
the full affine scan from fourth-power enumeration to third-power
multiset updates. The provided finite counters use fixed-width typed
arrays and ordinary exact-range integer arithmetic; they are test
programs for the documented sizes, not implementations for the
astronomical theorem onsets. The finite-search definition of
$a^{(3p)}$ permits arbitrary-precision integer arithmetic instead.

\section{Limits and unsuccessful routes}\label{sec:limits}
The bounds $f(n)\le0.44n^2$ eventually and $f(n)\ge0.31n^2$
eventually are not proved here. Nor do we prove that $f(n)/n^2$
converges. The following remarks describe obstructions encountered
by these particular arguments, not impossibility theorems for
stronger bounds.

\begin{remark}[Upper-bound obstructions]
The present table has $W_{\rm def}=0.002798309$, whereas a deficit
branch proving $0.44$ would require
$W_{\rm def}\ge B_0-0.44>0.006349540849362$.
Its deficit exceeds $0.003551231849362$, so these rows cannot prove
that target. Averaging inequalities of the form
$|\Sset|/n^2\le B_0-E$ and $|\Sset|/n^2\le B_0-A_x(e_x)$ adds no
constraint: a saving $s$ still requires the integrated inverse
threshold $\int v_s(x)\,dx\ge s$, where
$v_s(x)=\sup\{v:A_x(v)\ge s\}$.

High anchors alone are not mutually incompatible in the continuous
relaxation. The symmetric tent $a(t)=1-2|t-1/2|$ has nested maximal
cores and minimal exteriors. For symmetric windows of length $x$
and span $d$, exterior mass is $((1-x)^2+d^2)/2$ and common window
mass is $1/2-((1-x)^2+d^2)/2$; setting the latter to $1/4$ gives
$d=b(x)$. Thus the high-interval envelope can be attained in that
relaxation without information about low rows or distinct values.

Substituting an independent small diameter in the rank deficit and
a large diameter in the wing is unjustified, as is replacing
$\eta$ by the third root evaluated at $d_0$. Lemma~\ref{lem:coupled}
instead bounds $d-\eta$. A saturated wing $P=D_a$ can allow the gap
to extend outside the anchor union: \eqref{eq:coremass} excludes
containing the whole core but does not force a gap to meet the core.
This is why nonsaturation is proved before using
\eqref{eq:winggeometry}.

A floating scan using only the seed bound, at coefficient $0.4436$
and bin width $0.0002$, gave inverse-threshold integral about
$0.00265469$, below $B_0-0.4436\approx0.00274954$.
This is exploratory evidence, not a certified limitation.
Joint multipliers for two complements reproduced existing
single-complement or mass-only constraints; no further retained-wing
gain was established. A four-window floating relaxation used all
nonzero weight vectors in $\{0,1,2,3,4\}^4$, normalized by their sum,
gap increments $0.001$, start increments $0.003$, the parameter
triples $(x,v,\varepsilon)=(.33,.03,.03),(.31,.07,.03),
(.35,.015,.025),(.3,.09,.035)$, and
$y\in\{.01,.03,.05,.07,.1\}$. Off-center configurations survived;
for example the first triple at $y=.05$ allowed a tested gap $.1$
against a core-only root about $.07998$. Survival on a finite grid
is not a feasibility proof and supplies no symmetrization lemma.
Finally, an exploratory midpoint integration underestimated some
rectangle bounds; those provisional rows were discarded or reduced,
and only the passing directed-rounding certificate is used here.
\end{remark}

\begin{remark}[Lower-bound obstructions]
The coefficient $803/2592$ is below $31/100$ by exactly $13/64800$,
or $13/7200$ in units of $p^2$. Reducing the error term alone cannot
close this quadratic gap. The minimum central lift mass in the
present calculation is asymptotically $p^2/4$; no variance or
positive-fraction argument supplying an additional fixed gain is
proved here.

Fixing a translation in the four- and five-label estimates retains
all total-frequency planes without the controlled weight
\eqref{eq:weightnorm}. Deleting the translation average is therefore
not justified. The restricted interval average, together with the
estimates on every plane, is the proved replacement.
Optimizing finite affine scans does not establish a uniform
asymptotic gain, and unrestricted numerical maximizers are not
substituted for the theorem's parameter set.

Four- and five-row alternatives were also tested. For odd $p=2k+1$
and a permutation $\sigma$ of $1,\ldots,r$, take $z_1=x$,
$z_j=[2z_{j-1}]_p$ and the rows
\[
 (\sigma_1-1)p+x+1,\qquad
 (\sigma_j-1)p+k+1+z_j-z_{j-1}\ (2\le j<r),\qquad
 \sigma_rp-z_{r-1}.
\]
Since $[2z]_p-z$ permutes $[-k,k]$, these rows partition $1,\ldots,rp$
and have constant total $r(rp+1)/2$. All row orders at $r=4,5$,
$p=101,503$, $c=1$, $t=0,k$ give $576$ recorded counts; $18$ further
counts use selected orders at $p=1009,3331,5003$.
The order $(2,1,3,4)$ gives a finite ratio approximately $0.3125323102$
at $p=3331$ but $0.2895592821$ at $p=5003$. No uniform overlap
estimate for the additional prefix classes was proved, so this
route is neither established nor ruled out. Two interleaved triple
families and a second-moment approach were considered, but no
construction or estimate was completed for them.
\end{remark}

\section*{AI use and verification status}
The two proofs, the upper certificate and the accompanying programs
were developed with GPT-6 Astra (OpenAI). Each proof was reviewed in
isolation by two Claude (Anthropic) referees, both of whom returned
PASS. A separate Claude verification process returned PASS for the
upper certificate and PASS-WITH-REPAIRS for the lower result, with a
wording-only correction concerning the lower inequality's vacuity.
That correction is included in Theorem~\ref{thm:prime}'s discussion
and Section~\ref{sec:onsets}. The expository corrections from the
referees have also been incorporated. The review records report
independent exact certificate checks and finite combinatorial tests;
they do not constitute a formal proof verification. This manuscript
was prepared with GPT-6 from the reviewed proofs. Its presentation
has not undergone a separate referee round. No Lean formalization
has been carried out. The bounded literature search described in
Section~\ref{sec:related} does not certify priority.

\begin{thebibliography}{9}
\bibitem{Bloom}
T.~F. Bloom, \emph{Erd\H{o}s Problems}, Problem~34,
\url{https://www.erdosproblems.com/34}, accessed 12 October 2026
in the literature search described in Section~\ref{sec:related}.
\bibitem{Konieczny}
J. Konieczny, \emph{On consecutive sums in permutations},
Journal of Combinatorics \textbf{12} (2021), no.~3, 413--477.
Theorem~1.2, p.~415; Sections~1.5, 4.2 and 5.1.
\url{https://archive.intlpress.com/site/pub/files/_fulltext/journals/joc/2021/0012/0003/JOC-2021-0012-0003-a003.pdf}.
Preprint first posted 27 April 2015, \url{https://arxiv.org/abs/1504.07156}.
\bibitem{Hegyv}
N. Hegyv\'ari, \emph{On consecutive sums in sequences},
Acta Mathematica Hungarica \textbf{48} (1986), nos.~1--2, 193--200.
\url{https://doi.org/10.1007/BF01949064}.
\bibitem{Drakakis}
K. Drakakis, \emph{A review of Costas arrays},
Journal of Applied Mathematics \textbf{2006}, Article~26385,
Theorem~7.29, pp.~20--21.
\url{https://www.pure.ed.ac.uk/ws/portalfiles/portal/10369967/A_review_of_Costas_arrays.pdf}.
\bibitem{Golomb}
S.~W. Golomb, \emph{Algebraic constructions for Costas arrays},
Journal of Combinatorial Theory, Series~A \textbf{37} (1984), 13--21.
\bibitem{Dusart}
P. Dusart, \emph{Autour de la fonction qui compte le nombre de nombres
premiers}, doctoral thesis, Universit\'e de Limoges, 26 May 1998.
Th\'eor\`eme~1.9, printed p.~35 (supporting Proposition~1.10, p.~34).
\url{https://www.unilim.fr/pages_perso/pierre.dusart/Documents/T1998_01.pdf}.
\end{thebibliography}
\end{document}
